% Pista alg-23j-q4 · Àlgebra
% Procedència: PAU Catalunya 2023 · Convocatòria de juny · Sèrie 1 · Qüestió 4
\documentclass[11pt,a4paper]{article}
\usepackage{pista}
\pistainfo{Catalunya 2023}{Convocatòria de juny}{Sèrie 1}

\begin{document}

\section*{\textcolor{blauPista}{Qüestió 4}}

\noindent Sigui el sistema d'equacions lineals següent, que depèn del paràmetre real $\lambda$:
\[
\left.\begin{array}{r}
x + 2\lambda y + (2 + \lambda) z = 0 \\
(2 + \lambda) x + y + 2\lambda z = 3 \\
2\lambda x + (2 + \lambda) y + z = -3
\end{array}\right\}
\]

\subsection*{\textit{a)} \normalfont Discutiu el sistema per als diferents valors del paràmetre $\lambda$. \hfill\textnormal{[1,25~punts]}}

\pista{Escriu la matriu de coeficients $A$ i la matriu ampliada $A^{'}$. Calcula $\det(A)$ i troba els valors del paràmetre $\lambda$ que l'anul·len. Què pots dir del sistema quan $\det(A) \ne 0$? I per a cadascun dels valors que l'anul·len, compara els rangs d'$A$ i d'$A^{'}$.}

\subsection*{\textit{b)} \normalfont Per al cas $\lambda = -1$, resoleu el sistema, interpreteu-lo geomètricament i identifiqueu-ne la solució. \hfill\textnormal{[1,25~punts]}}

\pista{Fixa't que aquest apartat et demana tres coses ben diferents en una sola frase: «resoleu», «interpreteu geomètricament» i «identifiqueu-ne la solució». A l'apartat a) ja has vist de quin tipus és el sistema per a $\lambda = -1$: indica-ho. Si t'ha sortit compatible indeterminat, queda't amb dues equacions independents, pren una incògnita com a paràmetre (per exemple $z = t$) i expressa les altres dues en funció de $t$. Per a la interpretació geomètrica: cada equació del sistema representa un pla a l'espai, així que estàs estudiant la posició relativa de tres plans. Quina configuració correspon al cas que t'ha sortit (poden coincidir tots, tallar-se en una recta, no tenir cap punt comú\ldots)? Per identificar la solució, separa-la en la part que no depèn de $t$ i la que sí, $(x,y,z) = (\ldots) + t\,(\ldots)$, i interpreta què representa cada part.}

\end{document}
